\section{The Integers at the Bottom}

We have the tools: curved space, tilings, kaleidoscopes, addressing, graphs, automata, computation. Before we start building the vibe system on top of them, this chapter ties the geometry back to the deepest possible foundation, the plain whole numbers, and shows that the crystal is not a design choice but very nearly forced.

\subsection{The hunger for a non-arbitrary start}

Every theory of everything faces the same embarrassment. It starts with something, space, or particles, or a list of constants, and someone can always ask, why that and not something else? Why three forces? Why those masses? A truly satisfying foundation would start from something with nothing to choose, something so bare that it could not have been otherwise.

The barest thing there is, the thing with nothing to tune, is the whole numbers. One, two, three. You cannot pick different integers. They are not a design. So if a universe could be grown from the integers, with nothing added by hand at any step, that would be about as non-arbitrary as a foundation can get.

A word of care, because it is easy to misread. The claim is not that integers are the substance of reality. Integers do not feel anything, and the substance is feeling. The claim is that the integers are the right skeleton to describe the structure of the field of feeling, precisely because they smuggle in no assumptions. Experience is the flesh. The integers are the bones.

\subsection{The ladder}

Here is the climb, each rung adding exactly one idea, nothing arbitrary.

\begin{itemize}
\item \textbf{The integers.} The bedrock. No settings.
\item \textbf{Their symmetries.} The ways you can rearrange the integers without disturbing their structure form a clean, specific object, built from just two basic moves. We did not invent it. We asked the integers for their symmetries and this came back.
\item \textbf{From symmetries to mirrors.} Widen the idea of symmetry to include reflection, and you get a kaleidoscope, a Coxeter group, whose only settings are whole-number mirror angles.
\item \textbf{The kaleidoscope in curved space.} Let it run in hyperbolic space and the reflected copies of one tile fill the space with a crystal.
\item \textbf{The crystals.} Choosing the integers picks out specific crystals, the heptagrid, the dodecagrid. Among the most beautiful objects in mathematics, and any of them can serve as the fabric of space.
\end{itemize}

So the chain runs: integers, to their symmetries, to mirrors, to a hyperbolic crystal, to spacetime. Arithmetic becomes geometry. Nothing was tuned. You picked whole numbers, and a universe-shaped crystal fell out.

\subsection{The golden thread again}

Along the way the golden ratio keeps appearing, from independent directions, as the growth rate of the crystal, the most even spacing, and the spine of the addressing. Three separate roads, one number. That recurrence is the signature of a single ordered structure rather than a patchwork, and it is the kind of thing you hope to see when you claim a foundation is truly unified rather than assembled.

\subsection{Which integers carry the weight}

Not all integers matter equally. A few small ones are load-bearing, because a law needs only the smallest structure that can do its job. Three is the least number of values that can be ``against, neutral, toward,'' which is why the tone is three-valued. Seven is the least number of sides that curves the plane, which is why the first crystal is the heptagrid. Two, three, and four are the only dimensions a finite crystal can live in. The large integers are just addresses of cells, all equal. Importance is set by role, not by size, and the roles live among the small numbers.

\subsection{Where we stand}

That completes the tools. You now know what kind of space the crystal lives in and why, how it is generated from a few integers, how its cells are named and navigated, how it is really just a graph running a cellular automaton, and that such an automaton can compute anything. With all of that in hand, we can finally put feeling onto the crystal and watch the model itself come to life. That is Part III.

\textbf{Takeaway:} the crystal grows from the plain whole numbers, through their symmetries, to mirrors, to a hyperbolic crystal, with nothing tuned at any step and the golden ratio recurring as its signature. The integers are the bones, experience is the flesh, and the small integers carry the laws.
